\section*{अध्याय \ref{ch:b3:adaptive-immunity} --- अनुकूली प्रतिरक्षा और टीकाकरण}

\begin{solution}{exo:b3:adaptive-immunity:1}
एक \omterm{def:b3:adaptive-immunity:lymphocytes}{लसीकाणु}, एक विशिष्टता (नॉसल और लेडरबर्ग: हर कोशिका ने केवल एक ही
\omterm{def:b3:adaptive-immunity:lymphocytes}{प्रतिजन} के विरुद्ध \omterm{def:b3:adaptive-immunity:antibody}{प्रतिरक्षी} स्रवित किया); ग्राही \omterm{def:b3:adaptive-immunity:lymphocytes}{प्रतिजन} से मिलने से
पहले बन जाता है (टोनेगावा: वह जीन B~कोशिका में पुनर्विन्यस्त होता है,
\omterm{def:b3:adaptive-immunity:lymphocytes}{प्रतिजन} की अनुक्रिया में नहीं); \omterm{def:b3:adaptive-immunity:lymphocytes}{प्रतिजन} मेल खाते क्लोन चुनता है और
उन्हें कारक तथा \omterm{def:b3:adaptive-immunity:response}{स्मृति कोशिकाओं} में फैला देता है; स्व-क्रियाशील क्लोन
विकास के दौरान मिटा दिए जाते हैं।
\end{solution}

\begin{solution}{exo:b3:adaptive-immunity:2}
दो भारी और दो हल्की शृंखलाएँ, हर एक के सिरे पर कोई परिवर्ती प्रांत और
पीछे अचर प्रांत; छह अतिपरिवर्ती पाशों से बने दो एक जैसे बंधक स्थल; और
कोई Fc डंठल, जिसे दूसरी कोशिकाएँ तथा पूरक पढ़ते हैं। IgM: पहला
\omterm{def:b3:adaptive-immunity:antibody}{प्रतिरक्षी}, पंचलक, पूरक-सक्रियण। IgG: रक्त और ऊतक का मुख्य \omterm{def:b3:adaptive-immunity:antibody}{प्रतिरक्षी},
\omterm{def:b3:innate-immunity:complement}{ऑप्सोनीकरण}, उदासीनीकरण, अपरा पार करता है। IgA: द्विलक, जो श्लेष्मिक
पृष्ठों पर और दूध में स्रवित होता है। IgE: मास्ट कोशिकाओं से बँधा,
परजीवियों के विरुद्ध, एलर्जी का मध्यस्थ। IgD: अनुभवहीन B~कोशिकाओं पर,
जिसका ज्ञात प्रकार्य गौण है।
\end{solution}

\begin{solution}{exo:b3:adaptive-immunity:3}
वर्ग~I: सारी केंद्रकयुक्त कोशिकाएँ; कोशिकाद्रव्यी प्रोटीनों से
प्रोटिएसोम द्वारा काटे पेप्टाइड; 8--10 अवशेष; CD8 कोशिकाविषी
T~कोशिकाएँ पढ़ती हैं। वर्ग~II: \omterm{def:b3:innate-immunity:innate}{वृक्षाभ कोशिकाएँ}, \omterm{def:b3:innate-immunity:innate}{महाभक्षक}, B~कोशिकाएँ;
भीतर लेकर \omterm{def:b3:membrane-traffic:endocytosis}{अंतःकायों} में पचाए प्रोटीनों से पेप्टाइड; 13--25 अवशेष; CD4
सहायक T~कोशिकाएँ पढ़ती हैं।
\end{solution}

\begin{solution}{exo:b3:adaptive-immunity:4}
$R_{0}$ उन मामलों की औसत संख्या है जो कोई एक मामला पूरी तरह संवेद्य
समष्टि में पैदा करता है; देहली $p_{c} = 1 - 1/R_{0}$ है। गलसुआ
\qty{80}{\%}; कुकर खाँसी \qty{92}{\%}; 1918 का इन्फ़्लुएंज़ा \qty{60}{\%}।
\end{solution}

\begin{solution}{exo:b3:adaptive-immunity:5}
$50\times 20\times 5 = 5000$ भारी शृंखलाएँ, $60\times 4 = 240$ हल्की:
$1.2\times 10^{6}$ संयोजन; \omterm{def:b3:adaptive-immunity:antibody}{संधि-विविधता} के साथ $1.2\times
10^{10}$ --- यानी B~कोशिकाओं की संख्या से दस गुना, इसलिए वह प्राणी
अधिक से अधिक उसका दसवाँ भाग ही अभिव्यक्त करता है जो वह बना सकता था, और
लगभग हर कोशिका कोई अनन्य ग्राही ढोती है।
\end{solution}

\begin{solution}{exo:b3:adaptive-immunity:6}
$1000\times 10^{-3} = 1$ उत्परिवर्तन प्रति कोशिका प्रति विभाजन; बीस
विभाजनों में $20$। $10^{5}$ कोशिकाओं के बीच $2\times 10^{6}$
उत्परिवर्तन-घटनाएँ $3000$ संभव एकल प्रतिस्थापनों पर पड़ती हैं: हर एक
सैकड़ों बार प्रतिदर्शित होता है, और उसके अलावा दोहरे तथा तिहरे संयोजन
भी --- \omterm{def:b3:adaptive-immunity:germinal-centre}{जनन-केंद्र} उस ग्राही का पड़ोस पूरी तरह छान मारता है।
\end{solution}

\begin{solution}{exo:b3:adaptive-immunity:7}
T~कोशिकाओं का वरण \omterm{def:b3:adaptive-immunity:tolerance}{थाइमस} में इसलिए होता है कि वे परपोषी के अपने MHC
युग्मविकल्पियों पर पेप्टाइड पहचानें; किसी दूसरे व्यक्ति के
युग्मविकल्पियों की खाँचें भिन्न आकार की होती हैं, जो भिन्न पेप्टाइड
भिन्न दिशाओं में थामती हैं, इसलिए वही विषाणुजन्य पेप्टाइड दिखता ही
नहीं। किसी कलम के पराए MHC अणु स्वयं T~कोशिकाओं के बड़े अंश को ``किसी
अजीब पेप्टाइड के साथ स्व MHC'' जैसे दिखते हैं, इसलिए अस्वीकृति। वह
बहुरूपता इसलिए बनी रहती है कि अनेक युग्मविकल्पियों वाली कोई समष्टि किसी
भी रोगजनक के पेप्टाइडों का व्यापक प्रतिदर्श प्रस्तुत करती है और कोई
रोगजनक सबसे बच नहीं सकता; विषमयुग्मजी दुगुनी परास प्रस्तुत करते हैं और
उनका पक्ष लिया जाता है।
\end{solution}

\begin{solution}{exo:b3:adaptive-immunity:8}
(क) कोई \omterm{def:b3:adaptive-immunity:antibody}{V(D)J पुनर्संयोजन} नहीं: न B न T~कोशिकाएँ, यानी गंभीर संयुक्त
प्रतिरक्षा-न्यूनता। (ख) AIRE नहीं: \omterm{def:b3:adaptive-immunity:tolerance}{थाइमस} की कोशिकाएँ \omterm{def:b3:adaptive-immunity:lymphocytes}{ऊतक-प्रतिजन} नहीं
दिखा पातीं, स्व-क्रियाशील T~कोशिकाएँ बच निकलती हैं, अनेक अंगों के
विरुद्ध \omterm{def:b3:adaptive-immunity:tolerance}{स्वप्रतिरक्षा}। (ग) \omterm{def:b3:adaptive-immunity:tolerance}{FoxP3} नहीं: कोई नियामक T~कोशिकाएँ नहीं,
जल्दी शुरू होती घातक \omterm{def:b3:adaptive-immunity:tolerance}{स्वप्रतिरक्षा} और एलर्जी। (घ) AID नहीं: न
\omterm{def:b3:adaptive-immunity:antibody}{वर्ग-परिवर्तन}, न अतिउत्परिवर्तन --- भरपूर IgM, लगभग कोई IgG या IgA
नहीं, कम बंधुता के \omterm{def:b3:adaptive-immunity:antibody}{प्रतिरक्षी}, बार-बार लौटते संक्रमण। (ङ) CD4
T~कोशिकाएँ नहीं: B~कोशिकाओं या \omterm{def:b3:innate-immunity:innate}{महाभक्षकों} के लिए कोई मदद नहीं, कमज़ोर
प्रतिरक्षी-अनुक्रियाएँ और अवसरवादी संक्रमणों के प्रति संवेद्यता ---
यानी एड्स की प्रतिरक्षा-न्यूनता।
\end{solution}

\begin{solution}{exo:b3:adaptive-immunity:9}
$p_{c} = 1 - 1/6 = 0.83$; आच्छादन $0.83/0.9 = 93\,\%$। \qty{80}{\%}
आच्छादन पर प्रतिरक्षित अंश $0.72$ है, $R_{\text{प्रभावक}} = 6\times 0.28
= 1.7$। संवेद्यों के बीच वह महामारी ऐसे फैलती है मानो $R_{0}$ ही
$1.7$ हो (प्रतिरक्षित तो केवल संपर्क तनु करते हैं), और अंतिम-आकार
संबंध $s_{\infty} - \ln s_{\infty}/1.7 = 1$ से $s_{\infty}
\approx 0.31$ मिलता है: यानी संवेद्यों के कोई \qty{69}{\%}, और समष्टि
के \qty{19}{\%}, संक्रमित होते हैं।
\end{solution}

\begin{solution}{exo:b3:adaptive-immunity:10}
$\mathrm{d}s/\mathrm{d}t = -\beta si$ और $\mathrm{d}i/\mathrm{d}t
= \beta si - \gamma i$ से, $\mathrm{d}i/\mathrm{d}s = -1 + \gamma/(\beta s)$,
इसलिए $i + s - (\ln s)/R_{0}$ संरक्षित है; आरंभ पर $s = 1$, $i = 0$
और अंत पर $i = 0$, $s = s_{\infty}$ के साथ, $s_{\infty} - \ln
s_{\infty}/R_{0} = 1$। हल: $R_{0} = 1.5$, $s_{\infty} = 0.42$
(\qty{58}{\%} संक्रमित); $2$, $0.20$ (\qty{80}{\%}); $3$, $0.06$
(\qty{94}{\%}); $5$, $0.007$ (\qty{99.3}{\%})। सब नहीं: जैसे-जैसे
संवेद्य चुकते हैं, प्रभावी जनन संख्या कुछ के बचे रहते ही एक से नीचे
गिर जाती है, और महामारी उन तक पहुँचने से पहले बुझ जाती है।
\end{solution}

\begin{solution}{exo:b3:adaptive-immunity:11}
$1.5^{n} = 1000$: $n = \ln 1000/\ln 1.5 \approx 17$ क्रमिक सुधार।
प्रति कोशिका प्रति विभाजन एक उत्परिवर्तन और सौ में एक के सुधारक होने
के साथ, एक सुधार पाने के लिए प्रति चक्र कम से कम सौ कोशिकाएँ छाननी
पड़ती हैं; व्यवहार में हज़ारों, क्योंकि आधे उत्परिवर्तन बंधन नष्ट कर
देते हैं और वरण शोरभरा है। दो-दो विभाजनों के सत्रह चक्र कोई नौ दिन
लेते हैं --- यानी वही दो-तीन सप्ताह, जब अकुशलताएँ गिन ली जाएँ। चक्र
उत्परिवर्तन (गहरा क्षेत्र) को परीक्षण (हलका क्षेत्र) से अलग रखते हैं,
ताकि हर प्रभेद बढ़ने की छूट पाने से पहले \omterm{def:b3:adaptive-immunity:lymphocytes}{प्रतिजन} के सामने आँका जाए,
ठीक वैसे ही जैसे कोई आनुवंशिक कलन-विधि विचरण और वरण को बारी-बारी चलाती
है।
\end{solution}

\begin{solution}{exo:b3:adaptive-immunity:12}
स्त्रियाँ: कई प्रतिरक्षा-जीन X गुणसूत्र पर पड़े हैं, कुछ (TLR7)
\omterm{def:b3:genetic-engineering:transgenic}{निष्क्रियण} से बच जाते हैं और दोनों प्रतियों से अभिव्यक्त होते हैं, और
एस्ट्रोजन B~कोशिका का जीवित रहना बढ़ाते हैं --- जाँच: जिन पुरुषों में
एक अतिरिक्त X हो (क्लाइनफ़ेल्टर) उनमें ल्यूपस का ख़तरा स्त्रियों जैसा
होना चाहिए, और होता है। \omterm{def:b3:adaptive-immunity:mhc}{HLA}: संबद्ध युग्मविकल्पी संबंधित स्व पेप्टाइड
अच्छी तरह प्रस्तुत करते हैं, या \omterm{def:b3:adaptive-immunity:tolerance}{थाइमस} में उन्हें प्रस्तुत नहीं कर पाते
जिससे वे क्लोन मिटते नहीं --- जाँच: मानव युग्मविकल्पी के लिए
\omterm{def:b3:genetic-engineering:transgenic}{पारजीनी} चूहों में वह \omterm{def:b3:adaptive-immunity:lymphocytes}{प्रतिजन} देने पर वह रोग पनपता है। वह बढ़त: जल्दी
सूक्ष्मजीवी संसर्ग कम होना और बदले हुए \omterm{def:b3:microbiomes:symbiosis}{माइक्रोबायोम} (\omterm{def:b3:bacteriology:antibiotics}{प्रतिजैविक}, आहार,
शल्य प्रसव) कम नियामक T~कोशिकाएँ और अधिक \omterm{def:b3:innate-immunity:inflammation}{शोथ} देते हैं --- जाँचें:
मधुमेह-प्रवण चूहे रोगाणु-मुक्त या \omterm{def:b3:bacteriology:antibiotics}{प्रतिजैविकों} पर पाले जाएँ तो उनमें
मधुमेह अधिक होता है, और प्रवासियों के बच्चे एक ही पीढ़ी में नए देश की
घटना-दर अपना लेते हैं।
\end{solution}

\begin{solution}{pb:b3:adaptive-immunity:1}
\textbf{1.} $6000\times 320 = 1.9\times 10^{6}$; $\times 3\times 10^{4}
= 5.8\times 10^{10}$।
\textbf{2.} $10^{11}$ कोशिकाएँ, और संभव $5.8\times 10^{10}$: कोई
शरीर हर ग्राही लगभग दो बार रख सकता था, पर क्लोन असमान हैं और अधिकांश
अनुक्रम कभी बनते ही नहीं, इसलिए हर शरीर कोई बड़ा पर अधूरा प्रतिदर्श
अभिव्यक्त करता है --- और दो लोग अपने ग्राहियों का छोटा-सा अंश ही साझा
करते हैं।
\textbf{3.} $10^{11}/10^{5} = 10^{6}$ कोशिकाएँ; $10^{8}$ की किसी
ग्रंथि में लगभग $1000$।
\textbf{4.} $10^{9}/10^{5} = 10^{4}$ पूर्वज --- यानी सौ गुना कम;
नवजात के पुटक और \omterm{def:b3:innate-immunity:innate}{वृक्षाभ कोशिकाएँ} भी अपरिपक्व हैं, स्मृति कोई नहीं, और
T~कोशिका की मदद सीमित, इसलिए अनुक्रियाएँ धीमी तथा छोटी होती हैं, और
मातृ IgG वह खाई पाटता है।
\textbf{5.} वह V--D और D--J संधियों पर फैला है, जहाँ न्यूक्लियोटाइड
यादृच्छिक रूप से हटाए और जोड़े जाते हैं: उसका अनुक्रम जीनोम में कहीं
भी कूटबद्ध नहीं। बंधक स्थल के केंद्र में दोनों शृंखलाओं के बीच बैठा
होने से वही पाश एपिटोप के सबसे अधिक संपर्क में रहता है।
\textbf{6.} \omterm{def:b3:adaptive-immunity:lymphocytes}{प्रतिजन} से पहले विविधता चौड़ी होनी चाहिए ताकि जो भी आए उस
पर कुछ बैठे; \omterm{def:b3:adaptive-immunity:lymphocytes}{प्रतिजन} के बाद मेहनत उन थोड़े ग्राहियों को बदलने में
लगाना बेहतर है जिनका बाँधना पहले से पता है। पहले मोटी खोज, फिर बारीक
--- यह दोनों में से किसी एक अकेले से कहीं अधिक कुशल है।
\textbf{7.} सात दिन $21$ विभाजन हैं: $100\times 2^{21} = 2.1\times
10^{8}$ कोशिकाएँ; \qty{30}{\%}, यानी $6.3\times 10^{7}$ \omterm{def:b3:adaptive-immunity:response}{प्लाज़्मा
कोशिकाएँ}।
\textbf{8.} $6.3\times 10^{7}\times 2000\times \num{86400} = 1.1\times
10^{16}$ अणु प्रति दिन; $\times 1.5\times 10^{5}\times 1.66\times
10^{-24}$ g $= \qty{2.7}{mg}$; \qty{3}{L} में, प्रति दिन \qty{0.9}{mg/L}।
\textbf{9.} \qty{3}{L} में \qty{27}{mg}: \qty{9}{mg/L} $= \qty{0.009}{g/L}$,
यानी कुल IgG का हज़ारवाँ भाग --- किसी एक \omterm{def:b3:adaptive-immunity:lymphocytes}{प्रतिजन} का \omterm{def:b3:adaptive-immunity:antibody}{प्रतिरक्षी} पूरे का
छोटा-सा अंश होता है।
\textbf{10.} सौ गुना $\log_{2}100 = 6.6$ अर्ध-आयु हैं, लगभग
$140$ दिन। वर्षों का \omterm{def:b3:adaptive-immunity:antibody}{प्रतिरक्षी} मज्जा की दीर्घजीवी \omterm{def:b3:adaptive-immunity:response}{प्लाज़्मा कोशिकाओं}
से आता है, जो लगातार स्रवित करती हैं, और फिर संसर्ग पर दुबारा उद्दीपित
\omterm{def:b3:adaptive-immunity:response}{स्मृति कोशिकाओं} से।
\textbf{11.} नौ विभाजन: दिन 3 पर $10^{5}\times 512 = 5\times 10^{7}$
कोशिकाएँ --- यानी दिन 3 पर प्राथमिक की $5\times 10^{4}$ से हज़ार गुना,
और उसका चौथाई जहाँ प्राथमिक दिन 7 पर ही पहुँची थी।
\textbf{12.} पुनर्विन्यस्त VDJ C$_{\mu}$ के बगल में पड़ा है, इसलिए पहला
प्रतिलेख IgM है; C$_{\gamma}$ तक बदलने के लिए T~कोशिका की मदद और AID
चाहिए, जिनमें दिन लगते हैं। \omterm{def:b3:adaptive-immunity:response}{स्मृति कोशिकाएँ} पहले ही बदल चुकी हैं,
इसलिए उनकी संतति तुरंत IgG बनाती है।
\textbf{13.} $1000\times 10^{-3} = 1$ प्रति विभाजन; बीस में $20$।
\textbf{14.} सौ में एक; $10^{4}\times 0.01 = 100$ सुधरी कोशिकाएँ प्रति
विभाजन।
\textbf{15.} $1.5^{n} = 1000$: $n \approx 17$; $17\times \qty{12}{h}
\approx 8.5$ दिन।
\textbf{16.} आधी संततियाँ बाँधना खो देती हैं: वरण के बिना, बीस विभाजनों
के बाद उस वंश का $0.5^{20}
\approx 10^{-6}$ ही अब भी बाँधता। हलके क्षेत्र को हर चक्र पर हारे हुओं
को हटाना ही पड़ता है, ताकि केवल प्रकार्यशील और सुधरे ग्राही ही आगे बढ़ें।
\textbf{17.} पहली मात्रा ऊँची बंधुता और बदले हुए वर्ग की स्मृति
B~कोशिकाएँ छोड़ जाती है; दूसरी मात्रा उन्हीं से नए \omterm{def:b3:adaptive-immunity:germinal-centre}{जनन-केंद्र} बोती है,
और अतिउत्परिवर्तन तथा वरण किसी ऊँचे आरंभ-बिंदु से फिर चल पड़ते हैं,
इसलिए जो \omterm{def:b3:adaptive-immunity:antibody}{प्रतिरक्षी} निकलते हैं वे और भी बेहतर होते हैं।
\textbf{18.} कुछ संक्रमित लोग वर्षों के परिपक्वन के बाद अंततः उन
संरक्षित स्थलों के विरुद्ध \omterm{def:b3:adaptive-immunity:antibody}{प्रतिरक्षी} बना लेते हैं जिन्हें वह \omterm{def:b3:virology:virus}{विषाणु}
बदल नहीं सकता (व्यापक रूप से उदासीनकारी \omterm{def:b3:adaptive-immunity:antibody}{प्रतिरक्षी})। \omterm{def:b3:adaptive-immunity:germinal-centre}{जनन-केंद्र} को
दिशा दी जा सकती है: प्रतिरक्षाजनकों का कोई ऐसा क्रम, जो पहले सही
जनन-वंशीय पूर्वजों को जोड़े और फिर संरक्षित स्थल से बँधने को पुरस्कृत
करे, उस परिपक्वन को वर्षों के बजाय महीनों में उसी राह पर ले जा सकता है।
\textbf{19.} $p_{c} = 1 - 1/15 = 93.3\,\%$। एक मात्रा: $0.933/0.93 =
100.4\,\%$ --- अप्राप्य; दो मात्राएँ: $0.933/0.97 = 96.2\,\%$।
\textbf{20.} प्रतिरक्षित $0.90\times 0.97 = 0.873$; $R_{\text{प्रभावक}} = 15\times
0.127 = 1.9$: खसरा फैलता रहता है।
\textbf{21.} प्रति जन्म-समूह $60\,000\times (0.10 + 0.90\times 0.03) = 7600$
संवेद्य; पाँच वर्षों बाद $38\,000$ --- यानी कोई ऐसा समूह जिसमें किसी
आयातित मामले को अपने संपर्कों में शृंखलाएँ चलाने लायक संवेद्य मिल जाते
हैं।
\textbf{22.} संवेद्यों के बीच, $s_{\infty} - \ln s_{\infty}/1.9 =
1$ से $s_{\infty} \approx 0.23$ मिलता है: यानी उनके कोई \qty{77}{\%},
कोई $29\,000$ मामले।
\textbf{23.} प्रतिरक्षित $0.95\times 0.97 = 0.92$; $R_{\text{प्रभावक}} = 15\times
0.08 = 1.2$ --- फिर भी एक से ऊपर; देहली के लिए दो मात्राओं के साथ
\qty{96}{\%} चाहिए, और खसरा वही रोग है जो आख़िरी कुछ प्रतिशत की सज़ा
देता है।
\textbf{24.} शिशु केवल अपने आसपास के लोगों की प्रतिरक्षा से सुरक्षित
हैं (और कुछ महीनों तक मातृ \omterm{def:b3:adaptive-immunity:antibody}{प्रतिरक्षी} से): यदि संचरण की शृंखलाएँ बन ही
न सकें, तो कोई मामला उन तक पहुँचता नहीं। \qty{85}{\%} आच्छादन पर
प्रतिरक्षित अंश $0.82$ है और $R_{\text{प्रभावक}} = 2.6$: प्रकोप चलते हैं,
और शिशु, अटीकाकृत तथा सबसे नाज़ुक, संक्रमित होते हैं।
\textbf{25.} संग्रह लगभग $6\times 10^{10}$; किसी एक अनुक्रिया की
\omterm{def:b3:adaptive-immunity:response}{प्लाज़्मा कोशिकाओं} से दिन भर में \qty{2.7}{mg} IgG; और \qty{96}{\%}
दो-मात्रा आच्छादन।
\end{solution}
